\ifdefined\CRevisionRound\else\def\CRevisionRound{2}\fi
\documentclass[10pt]{article}
\usepackage[margin=0.62in]{geometry}
\usepackage{amsmath,amssymb,amsthm,microtype,titlesec}
\usepackage[T1]{fontenc}
\usepackage{lmodern}
\pdfvariable suppressoptionalinfo 512\relax
\pdfvariable trailerid {[<c326c326c326c326c326c326c326c326><c326c326c326c326c326c326c326c326>]}
\newtheorem{theorem}{Theorem}
\newtheorem{lemma}{Lemma}
\newcommand{\Prob}{\mathbb P}
\newcommand{\norm}[1]{\lVert#1\rVert}
\titlespacing*{\section}{0pt}{1.6ex plus .4ex}{.7ex}
\title{The Two-Site Inclusion Process:\\Complete Hahn Spectrum and the Absorbing Face}
\author{HCS-C326 theorem-and-evidence package}
\date{3 September 2026 --- revision round \CRevisionRound}
\begin{document}
\maketitle
\begin{abstract}
We diagonalize the convention-locked two-site symmetric inclusion process.
The result includes its beta-binomial law, every Hahn mode and finite-sum norm,
the full semigroup kernel, sharp $L^2$ decay, and the singular zero-parameter
absorbing limit. Exact rational grids audit, but do not prove, the theorem.
\end{abstract}

\section{Frozen generator and reversible law}
Fix $N\in\mathbb Z_{\geq0}$ and $\alpha>0$. State $x\in S_N=\{0,\ldots,N\}$
is the first-site occupancy; the second is $N-x$. Our unscaled continuous-time
generator is
\begin{equation}\label{eq:L}
 (Lf)(x)=b_x[f(x+1)-f(x)]+d_x[f(x-1)-f(x)],\quad
 b_x=(N-x)(\alpha+x),\quad d_x=x(\alpha+N-x),
\end{equation}
with unavailable endpoint terms zero. Put
\begin{equation}\label{eq:pi}
 \pi(x)=\frac{N!(\alpha)_x(\alpha)_{N-x}}
 {x!(N-x)!(2\alpha)_N}.
\end{equation}

\begin{theorem}\label{thm:main}
For $N\geq1$, \eqref{eq:pi} is the unique reversible law. Define
\[
 H_j(x)={}_3F_2\!(\begin{smallmatrix}-j,j+2\alpha-1,-x\\
                         \alpha,-N\end{smallmatrix};1),\quad
 h_j=\sum_{x=0}^N\pi(x)H_j(x)^2,\quad 0\leq j\leq N.
\]
Then $h_j>0$, the $H_j$ are a complete orthogonal basis, and
\[
 LH_j=-\lambda_jH_j,\qquad \lambda_j=j(j-1+2\alpha).
\]
The complete transition kernel is
\begin{equation}\label{eq:kernel}
 p_t(x,y)=\pi(y)\sum_{j=0}^Ne^{-\lambda_jt}
 \frac{H_j(x)H_j(y)}{h_j}.
\end{equation}
For centered $f$, $\norm{e^{tL}f}_{2,\pi}\leq e^{-2\alpha t}\norm f_{2,\pi}$;
the exponent is sharp. At $\alpha=0$, the endpoints absorb, absorption is
almost sure, and $\Prob_x(X_\infty=N)=x/N$ for $N\geq1$. Moreover
$\{c\delta_0+(1-c)\delta_N:0\leq c\leq1\}$ is the complete stationary family,
and
$\pi_{\alpha,N}\Rightarrow(\delta_0+\delta_N)/2$ as $\alpha\downarrow0$.
For $N=0$ the process and law are the singleton.
\end{theorem}

The adjacent identity $\pi(x)b_x=\pi(x+1)d_{x+1}$ follows by cancelling
Pochhammer factors; Chu--Vandermonde gives the denominator in \eqref{eq:pi}.
For $N\geq1$ all interior edge rates are positive, proving irreducibility,
uniqueness, and self-adjointness in $L^2(\pi)$.

The evidence grid uses $\alpha\in\{1/2,1,3/2,2\}$ and $0\leq N\leq8$:
36 parameter rows, 180 states, and 180 full Hahn vectors. All entries are
exact rationals independently reconstructed; finite evidence is not a proof.

\ifnum\CRevisionRound>0
\section{Full Hahn diagonalization}
\begin{lemma}\label{lem:tri}
The degree-$j$ polynomial filtration is invariant under $L$, and the diagonal
entry on its one-dimensional quotient is $-j(j-1+2\alpha)$.
The terminating series defining $H_j$ has degree $j$, satisfies
$LH_j=-j(j-1+2\alpha)H_j$, and obeys $H_j(0)=1$.
\end{lemma}
\begin{proof}
For $f(x)=x^j$, the prospective $x^{j+1}$ terms in \eqref{eq:L} cancel. The
$x^j$ coefficient is
\[
 j[(N-\alpha)-(N+\alpha)]-j(j-1)=-j(j-1+2\alpha).
\]
For identification, write the finite expression
\[
 H_j(x)=\sum_{k=0}^j
 \frac{(-j)_k(j+2\alpha-1)_k(-x)_k}
      {(\alpha)_k(-N)_k k!}.
\]
Put $\phi_k(x)=(-x)_k$. Direct forward and backward differencing gives
\[
 L\phi_k=-k(k-1+2\alpha)\phi_k
 -k(N-k+1)(\alpha+k-1)\phi_{k-1}.
\]
For $H_j=\sum_kc_k\phi_k$, the eigenvalue equation requires
\[
 \frac{c_{k+1}}{c_k}=
 \frac{\lambda_j-\lambda_k}{(k+1)(N-k)(\alpha+k)}
 =\frac{(j-k)(j+k+2\alpha-1)}{(k+1)(N-k)(\alpha+k)}.
\]
This is exactly the ratio of consecutive coefficients in the displayed
terminating sum, by $(a)_{k+1}=(a+k)(a)_k$. Consequently every coefficient
on the left of
\[
 b_x[H_j(x+1)-H_j(x)]+d_x[H_j(x-1)-H_j(x)]
 +j(j-1+2\alpha)H_j(x)
\]
is zero. This proves the finite difference identity. Its denominators do not
vanish for $j\leq N$ and $\alpha>0$; its final coefficient is nonzero, while
substitution $x=0$ gives one.
\end{proof}

Because $\lambda_{j+1}-\lambda_j=2(j+\alpha)>0$, these are all the distinct
diagonal values from Lemma~\ref{lem:tri}. Self-adjointness makes their
eigenvectors orthogonal. There are $N+1$ of them, hence they are complete, and
each displayed finite-sum norm $h_j$ is positive. Expanding any function in
this normalized basis proves \eqref{eq:kernel}. Parseval on the constants'
orthogonal complement gives the stated bound because $\lambda_1=2\alpha$;
$H_1(x)=1-2x/N$ gives equality and therefore sharpness. When $N=1$, this says
exactly that the spectrum is $\{0,2\alpha\}$ and $\pi=(1/2,1/2)$.

\section*{Revision certificate: complete Hahn spectrum and semigroup}
This round adds polynomial triangularity, the terminating difference identity,
orthogonality with positive finite-sum norms, the full kernel, and sharp decay.
\fi

\ifnum\CRevisionRound>1
\section{The singular face, ownership, and Route A}
At $\alpha=0$, $b_x=d_x=x(N-x)$. Thus both endpoints absorb and, from an
interior state, the embedded chain is the finite simple symmetric walk.
Absorption is almost sure. Since $Lx=0$, bounded optional stopping yields
$x=N\Prob_x(X_\infty=N)$. A stationary law cannot charge the transient
interior: stationarity makes that mass time-independent, whereas finite-state
almost-sure absorption makes it tend to zero. Thus all and only stationary
laws are $c\delta_0+(1-c)\delta_N$, $0\leq c\leq1$. For fixed $N\geq1$, each endpoint numerator in
\eqref{eq:pi} is asymptotic to $\alpha/N$, the denominator to $2\alpha/N$,
and every interior numerator is $O(\alpha^2)$; this proves the weak limit.
For $N=0$ only $\delta_0$ exists. At $N=1$ the hitting formula remains valid.

The inclusion model is sourced to Giardina--Redig--Vafayi~\cite{grv}; the Hahn
normalization, representation, and difference equation are checked against
DLMF Sections 18.19, 18.20.5, and 18.22(ii)~\cite{dlmf}. We claim no priority. C253
owns Moran fixation, C263 a growing Polya urn, C285 Gordon--Newell product form,
and C322 Kac sphere collisions; none owns this theorem. We assert no multisite,
open-boundary, or condensation-scaling result.

The Route-A tuple is $(\mathrm{FAIL},\mathrm{FAIL},\mathrm{FAIL},
\mathrm{FAIL},\mathrm{FORMAL\ HINT})$, with verdict
\texttt{ROUTE\_A\_REJECTED}; Route B is false. The final entry records only a
finite self-adjoint spectral analogy, not a Hilbert--Polya operator. Under
\texttt{NO\_BAD\_EULER\_\allowbreak OR\_ROOT\_NUMBER}, no target local datum, Euler factor,
root number, automorphy, divisor, functional equation, or zero match is asserted.

\section*{Revision certificate: absorbing face and scope closure}
This final round proves absorption, hitting probabilities, the stationary weak
limit, all degenerate cases, source and collision ownership, and the firewall.
\fi

\ifnum\CRevisionRound=0
\section*{Revision certificate: frozen inclusion chain and reversible law}
This original round fixes the generator and clock, derives beta-binomial
detailed balance, states the complete theorem, and records the exact grid.
\fi

\begin{thebibliography}{2}
\bibitem{grv} C. Giardina, F. Redig, and K. Vafayi, Correlation inequalities
for interacting particle systems with duality, \emph{J. Stat. Phys.} 141
(2010); DOI: 10.1007/s10955-010-0055-0; arXiv:0906.4664.
\bibitem{dlmf} NIST Digital Library of Mathematical Functions, Sections
18.19, 18.20.5, and 18.22(ii), Hahn definitions, representation, and
difference equations.
\end{thebibliography}
\end{document}
